\section{Presentations: extensive form and Kuhn realisation}
\label{sec:presentations}

The kernel games of \S\ref{sec:fosg-strategies} are the primary
strategic-form interface.  Two further results are needed to make
the construction interoperate with the textbook tradition: an
extensive-form presentation of the FOSG view, and a finite Kuhn
realisation theorem that links mixed and behavioural strategies.
The native strategic layer is not defined through these presentations;
their connection to native strategies is an explicit bridge obligation.

\subsection{The canonical EFG presentation}
\label{sec:pres:efg}

For each checked program $g$, the codebase produces a canonical
extensive-form game $\mathsf{efg}(g)$ by serialising the bounded
event-graph FOSG view: each FOSG round becomes a block of $|
\Players|$ sequential decision opportunities (each over the player's
local round-action alphabet, augmented with $\mathsf{none}$ for
inactive players) followed by a chance node implementing the FOSG
transition.

\begin{theorem}[EFG presentation adequacy]
\label{thm:efg}
For every checked program $g$ and every bounded event-graph FOSG
behavioural profile $\beta$, the translated EFG outcome kernel
projected to the payoff-vector outcome agrees with the bounded FOSG
outcome kernel:
\[
  \mathrm{map}(\mathrm{publicOf})
    \,\bigl(\mathrm{outcomeKernel}_{\mathsf{efg}(g)}\,\beta\bigr)
  \;=\;
  K^{\mathrm{FOSG}}_{\mathrm{beh}}(g)(\beta),
\]
where $\beta$ is first translated by the canonical FOSG-to-EFG profile
translation.
\end{theorem}

The proof is by induction on the bounded FOSG depth, exploiting the
fact that translated legal FOSG behavioural strategies put all mass
on $\mathsf{none}$ for inactive players, so the administrative
decision opportunities at the EFG layer add no strategic freedom.

\designnote{The EFG bridge is intentionally a \emph{view}, not the
semantic owner.  The authoritative semantics is the event-graph
machine; the EFG is a presentation that respects the payoff-vector
outcome law of the bounded FOSG presentation.  Relating that FOSG law
to the native Vegas strategy carrier is the same native/FOSG bridge
obligation needed by the native Kuhn statement.  This preserves the two
strategically irrelevant features of the EFG — the arbitrary
linearisation of independent same-round events, and the introduction of
administrative no-op decision opportunities — without making them part
of the underlying game.}

\subsection{Mixed extension and the finite Kuhn theorem}
\label{sec:pres:kuhn}

The \emph{mixed extension} of a finite kernel game $\KG$ replaces
each strategy space $S_p$ by $\PMF{S_p}$ and pushes the outcome
kernel through the independent product:
\begin{equation}
\label{eq:mixed-extension}
  \Kmix(\mu)
  \;=\;
  \bigl({\textstyle\bigotimes_p}\,\mu_p\bigr) \bind \Kpure.
\end{equation}

Kuhn's classical theorem~\citep{Kuhn1953} says that in finite
extensive-form games with perfect recall, mixed strategies over pure
contingent plans and behavioural strategies are realisationally
equivalent.  The unconditional theorem currently mechanised in Vegas is
the bounded event-graph FOSG counterpart.  The native public theorem is
stated over the native kernel games and takes the native/FOSG
strategy-equivalence bridge as an explicit argument.

\begin{theorem}[Finite native Kuhn realisation, conditional bridge]
\label{thm:kuhn}
For every checked program $g$, assuming the native/FOSG strategy-equivalence
bridge for $g$, every independent mixed profile
$\mu \in \prod_p \PMF{S_p^{\mathrm{pure}}}$ over pure strategies of
$\mathrm{pureKernelGame}(g)$ admits a behavioural profile $\beta$ of
$\mathrm{pmfBehavioralKernelGame}(g)$ with the same payoff-vector outcome
distribution:
\[
  \Kbeh(\beta)
  \;=\;
  \bigl({\textstyle\bigotimes_p}\,\mu_p\bigr) \bind \Kpure.
\]
\end{theorem}

\begin{proof}[Proof sketch]
On the bounded FOSG carrier, $\beta$ is constructed pointwise on
reachable information states: at each
$I \in \mathrm{ReachableInfoState}(p)$, $\beta_p(I)$ is the posterior
distribution over actions induced by $\mu_p$ conditional on the
history class realising $I$.  The FOSG proof uses the
\emph{LegalObservable} theorem for checked event-graph FOSGs.  The
native theorem then applies the explicit bridge obligation, which is
the still-separate proof that native round-view strategies and bounded
FOSG strategies realise the same mixed and behavioural outcome laws.
\end{proof}

\Cref{thm:kuhn} is the intended contract between the pure-strategy API
and the PMF-behavioural analysis API.  The current Lean statement makes
the remaining native/FOSG bridge visible rather than hiding it in prose.
It does not by itself construct an equilibrium; existence follows from
composing realisation with a fixed-point theorem
(\S\ref{sec:fosg-strategies}).

\paragraph{Lean correspondence.}
Supplement~\S6.4 (canonical EFG presentation), \S7.6 (mixed extension
and Kuhn realisation: \TT{kuhn\_finite}, \TT{NativeFOSGKuhnBridge},
\TT{eventGraphEFGAt\_outcomeKernel\_map\_publicOutcome}).
